\noindent\textbf{Part 1}
\begin{description}
\item[(a)] The spiral arm spacing of AFGL 3068 is
  \[ \Delta r_{\mathrm{arm}}[\mathrm{au}]
     = \Delta\theta[\mathrm{arcsec}] \times d[\mathrm{pc}]
     = 2.29 \times 1000 = 2290 \]
  \hfill(1 point)

  When the angle $\phi$ increases by $2\pi$, the radius $r$ increases by an
  amount equal to the arm spacing. Therefore, we have
  \[ b = \frac{\Delta r_{\mathrm{arm}}}{2\pi} = \frac{2290}{2\pi}
     \approx 364\,\mathrm{au\,rad^{-1}} \]
  \hfill(1 point)

\item[(b)] Consider a shell element located on the orbital plane. Because it
  moves radially outward with velocity $V_w$, its radial distance at a time
  $t$ is
  \[ r(t) = r_p + V_w t \]
  \hfill(1 point)

  At the same time, the red giant rotates an angle
  \[ \phi(t) = \frac{V_p}{r_p}\,t \]
  \hfill(1 point)

  Extracting $t$ from this equation and substituting it into the previous
  equation yields
  \[ r(t) = r_p + \frac{r_p V_w}{V_p}\,\phi(t) \]
  \hfill(1 point)

  On the orbital plane, shell elements all obey an equation of the form given
  above, but with different values of $\phi$, as they are ejected at different
  times. At the observation time, these elements trace an Archimedean spiral
  described by
  \[ r = r_p + \frac{r_p V_w}{V_p}\,\phi \]
  \hfill(1 point)
  \[ \Longrightarrow a = r_p, \qquad b = \frac{r_p V_w}{V_p} \]

\item[(c)] From b), the constant $b$ is
  \[ b = \frac{r_p V_w}{V_p} \]
  The orbital period of the red giant is given by
  \[ T = \frac{2\pi r_p}{V_p} = \frac{2\pi b}{V_w} \]
  \hfill(1 point)

  Substituting $V_w = 14.0\,\mathrm{km\,s^{-1}} \approx
  2.95\,\mathrm{au\,yr^{-1}}$ and $b = 364\,\mathrm{au\,rad^{-1}}$ into the
  above formula yields
  \[ T = \frac{2\pi \times 364}{2.95} \approx 775\,\mathrm{yr} \]
  \hfill(1 point)

\item[(d)] Applying Kepler's third law to the binary system, the separation
  between the two stars is given by
  \[ d^3 = M_{\mathrm{tot}} T^2 \]
  \hfill(1 point)
  \[ \Rightarrow d = \sqrt[3]{M_{\mathrm{tot}} T^2}
     = \sqrt[3]{6.6 \times 775^2} \approx 158\,\mathrm{au} \]
  \hfill(1 point)
\end{description}

\noindent\textbf{Part 2}
\begin{description}
\item[(a)] The coordinates of a point on the shell can be written as
  \[ x = r\sin\theta\cos\phi = \frac{r_p V_w}{V_p}\,\phi\sin\theta\cos\phi \]
  \hfill(1.5 points)
  \[ y = r\sin\theta\sin\phi = \frac{r_p V_w}{V_p}\,\phi\sin\theta\sin\phi \]
  \hfill(1.5 points)
  \[ z = r\cos\theta = \frac{r_p V_w}{V_p}\,\phi\cos\theta \]
  \hfill(1.5 points)

  Because all shell elements move radially outward with velocity $V_w$, the
  components of the velocity in the Cartesian coordinates are given by
  \[ V_x = V_w \sin\theta\cos\phi \]
  \hfill(0.5 points)
  \[ V_y = V_w \sin\theta\sin\phi \]
  \hfill(0.5 points)
  \[ V_z = V_w \cos\theta \]
  \hfill(0.5 points)

\item[(b)] Choose the sky plane coordinates $(X, Y)$ with the $X$-axis
  identical to the $x$-axis. The line-of-sight velocity is
  \[ V_{\mathrm{los}} = V_y \sin i - V_z \cos i
     = V_w(\sin\theta\sin\phi\sin i - \cos\theta\cos i) \]
  \hfill(2 points)

  From $V_{\mathrm{los}} = 0$, we have
  \[ \sin\theta\sin\phi\sin i = \cos\theta\cos i \]
  \hfill(1 point)

  The sky plane coordinates $(X, Y)$ of a shell element are given by
  \[ X = x = r\sin\theta\cos\phi \]
  \hfill(1 point)
  \[ Y = y\cos i + z\sin i = r(\sin\theta\sin\phi\cos i + \cos\theta\sin i) \]
  \hfill(1 point)

  Substituting $\cos\theta = \sin\theta\sin\phi\tan i$ into $Y$ coordinate,
  \[ Y = r\sin\theta\sin\phi\,(\cos i + \tan i \sin i)
     = \frac{r\sin\theta\sin\phi}{\cos i} \]
  \hfill(1 point)

  The sky plane coordinates $(X, Y)$ satisfy
  \[ \frac{X^2}{(r\sin\theta)^2}
     + \frac{Y^2}{\left(\dfrac{r\sin\theta}{\cos i}\right)^2} = 1 \]
  \hfill(1 point)

\item[(c)] From b), the ratio of the semi-major axis and the semi-minor axis
  is given by
  \[ \frac{1}{\cos i} \]
  \hfill(1 point)

  For the spiral pattern of AFGL 3068, we have
  \[ \frac{1}{\cos i} = 1.1 \Rightarrow i \approx 24.6^\circ \]
  \hfill(1 point)
\end{description}
